\documentclass[12pt,a4paper]{article}
\usepackage[utf8]{inputenc}
\usepackage[T1]{fontenc}
\usepackage{geometry}
\geometry{margin=2.5cm}
\usepackage{hyperref}
\usepackage{enumitem}
\usepackage{fancyhdr}
\usepackage{titlesec}
\usepackage{array}
\usepackage{longtable}

\title{\textbf{User Stories Détaillées - Système MineTrack}}
\author{Documentation du Projet}
\date{\today}

\pagestyle{fancy}
\fancyhf{}
\rhead{MineTrack}
\lhead{User Stories}
\rfoot{Page \thepage}

\titleformat{\section}{\normalfont\Large\bfseries}{\thesection}{1em}{}[\titlerule]

\begin{document}



\section{Introduction}
Ce document présente les \textit{User Stories} détaillées pour le système de gestion minière \textbf{MineTrack}. Le système se compose d'une application Web pour l'administration et le suivi global, et d'une application Mobile pour les opérations sur le terrain.

\section{Acteurs du Système}
\begin{itemize}[label=\textbullet]
    \item \textbf{Administrateur :} Gère les sites, le personnel et les configurations globales (Web).
    \item \textbf{Chef d'Équipe :} Supervise les opérations quotidiennes, le pointage et la production (Web \& Mobile).
    \item \textbf{Ouvrier :} Consulte son planning et effectue son pointage (Mobile).
\end{itemize}

\section{User Stories - Application Web (Administration)}

\begin{longtable}{|p{3cm}|p{12cm}|}
\hline
\textbf{ID} & \textbf{US-W01 : Gestion des Sites} \\ \hline
\textbf{En tant que} & Administrateur \\ \hline
\textbf{Je veux} & Créer, modifier et supprimer des sites miniers \\ \hline
\textbf{Critères d'acceptation} & 
\begin{itemize}[nosep, leftmargin=*]
    \item Pouvoir saisir le nom, le code unique, le minerai et les coordonnées.
    \item Définir la capacité et le nombre de travailleurs.
    \item Voir le statut (actif/inactif) de chaque site.
\end{itemize} \\ \hline
\end{longtable}

\begin{longtable}{|p{3cm}|p{12cm}|}
\hline
\textbf{ID} & \textbf{US-W02 : Dashboard Global} \\ \hline
\textbf{En tant que} & Administrateur / Chef d'Équipe \\ \hline
\textbf{Je veux} & Visualiser les statistiques en temps réel de tous les sites \\ \hline
\textbf{Critères d'acceptation} & 
\begin{itemize}[nosep, leftmargin=*]
    \item Affichage de la production totale vs objectif.
    \item Graphiques d'évolution mensuelle.
    \item Alerte visuelle pour les incidents critiques.
\end{itemize} \\ \hline
\end{longtable}

\begin{longtable}{|p{3cm}|p{12cm}|}
\hline
\textbf{ID} & \textbf{US-W03 : Gestion du Personnel} \\ \hline
\textbf{En tant que} & Administrateur \\ \hline
\textbf{Je veux} & Gérer les dossiers des employés et leurs affectations \\ \hline
\textbf{Critères d'acceptation} & 
\begin{itemize}[nosep, leftmargin=*]
    \item Enregistrer les informations personnelles et professionnelles.
    \item Affecter un employé à un site et à une équipe spécifique.
    \item Suivre la date d'embauche et le statut (actif/en congé).
\end{itemize} \\ \hline
\end{longtable}

\begin{longtable}{|p{3cm}|p{12cm}|}
\hline
\textbf{ID} & \textbf{US-W04 : Suivi de la Production} \\ \hline
\textbf{En tant que} & Administrateur \\ \hline
\textbf{Je veux} & Suivre et analyser les données de production de tous les sites \\ \hline
\textbf{Critères d'acceptation} & 
\begin{itemize}[nosep, leftmargin=*]
    \item Visualiser les tonnages extraits et le nombre de camions par site et par shift.
    \item Comparer la production réelle par rapport aux objectifs fixés.
    \item Filtrer les données de production par période (jour, semaine, mois).
    \item Exporter les données de production sous format Excel ou PDF.
\end{itemize} \\ \hline
\end{longtable}

\begin{longtable}{|p{3cm}|p{12cm}|}
\hline
\textbf{ID} & \textbf{US-W05 : Gestion des Rapports d'Incidents} \\ \hline
\textbf{En tant que} & Administrateur \\ \hline
\textbf{Je veux} & Gérer et traiter les rapports d'incidents envoyés depuis le terrain \\ \hline
\textbf{Critères d'acceptation} & 
\begin{itemize}[nosep, leftmargin=*]
    \item Centraliser tous les rapports d'incidents avec leur niveau de gravité.
    \item Changer le statut d'un incident (Ouvert, En cours, Résolu, Clôturé).
    \item Ajouter des commentaires ou des instructions de suivi sur un rapport.
    \item Générer des rapports de sécurité périodiques.
\end{itemize} \\ \hline
\end{longtable}

\section{User Stories - Application Mobile (Opérations)}

\begin{longtable}{|p{3cm}|p{12cm}|}
\hline
\textbf{ID} & \textbf{US-M01 : Pointage de l'Équipe par le Chef} \\ \hline
\textbf{En tant que} & Chef d'Équipe \\ \hline
\textbf{Je veux} & Valider la présence de tous les membres de mon équipe au début du shift \\ \hline
\textbf{Critères d'acceptation} & 
\begin{itemize}[nosep, leftmargin=*]
    \item Accéder à la liste complète des membres de l'équipe assignés au shift actuel.
    \item Effectuer le pointage manuel pour chaque ouvrier (Présent, Absent, Retard).
    \item Afficher les employés en repos (OFF) pour permettre leur rappel en cas de besoin.
    \item Horodatage automatique de la validation effectuée par le Chef d'Équipe.
    \item Possibilité de modifier le statut en cours de shift (ex: départ anticipé ou pause).
\end{itemize} \\ \hline
\end{longtable}

\begin{longtable}{|p{3cm}|p{12cm}|}
\hline
\textbf{ID} & \textbf{US-M02 : Saisie de Production} \\ \hline
\textbf{En tant que} & Chef d'Équipe \\ \hline
\textbf{Je veux} & Déclarer les quantités extraites durant mon shift \\ \hline
\textbf{Critères d'acceptation} & 
\begin{itemize}[nosep, leftmargin=*]
    \item Sélectionner la zone de travail.
    \item Saisir le tonnage et le nombre de camions chargés.
    \item Consulter l'historique des saisies de la journée.
\end{itemize} \\ \hline
\end{longtable}

\begin{longtable}{|p{3cm}|p{12cm}|}
\hline
\textbf{ID} & \textbf{US-M03 : Déclaration d'Incident} \\ \hline
\textbf{En tant que} & Chef d'Équipe / Ouvrier \\ \hline
\textbf{Je veux} & Signaler immédiatement un incident de sécurité ou technique \\ \hline
\textbf{Critères d'acceptation} & 
\begin{itemize}[nosep, leftmargin=*]
    \item Choisir le type d'incident et définir le niveau de gravité selon l'échelle standard (Faible, Moyen, Élevé, Critique).
    \item Saisir une description détaillée, la localisation précise et les éventuels dommages.
    \item Associer une couleur visuelle au rapport en fonction de la gravité (Vert, Jaune, Orange, Rouge).
    \item Envoi instantané d'une alerte prioritaire au dashboard web.
\end{itemize} \\ \hline
\end{longtable}

\begin{longtable}{|p{3cm}|p{12cm}|}
\hline
\textbf{ID} & \textbf{US-M04 : Consultation du Planning} \\ \hline
\textbf{En tant que} & Ouvrier \\ \hline
\textbf{Je veux} & Consulter mon planning hebdomadaire \\ \hline
\textbf{Critères d'acceptation} & 
\begin{itemize}[nosep, leftmargin=*]
    \item Voir le shift assigné pour chaque jour de la semaine.
    \item Connaître les équipements assignés et la zone de travail.
    \item Recevoir des notifications en cas de changement.
\end{itemize} \\ \hline
\end{longtable}

\end{document}
